\documentclass[11pt]{article}
\usepackage{mathblatt}
\def\mitzone{1}
\begin{document}
\blattkopf{Quadratische Funktionen}{Gesamt}
\weit
\verzeichniszeile{\verz{zone}{Kennst du schon (Nr.~1–10)} \verztrenn \verz{e1}{1 Normalparabel und Streckfaktor (Nr.~11–12)} \verztrenn \verz{e2}{2 Scheitelpunktform (Nr.~13–14)} \verztrenn \verz{e3}{3 Normalform (Nr.~15–16)} \verztrenn \verz{e4}{4 Nullstellen und Schnittpunkte berechnen (Nr.~17–18)} \verztrenn \verz{abhaken}{Das kann ich}}
% Zone „Das kennst du schon“ – aus bank/quadratische-funktionen/zone.jsonl (zusammenbau v0.1)
\einheitenkopf[zone][Kennst du schon]{Das kennst du schon}
\setcounter{aufgabe}{0}

%% TODO Titel: Ich-kann-Satz fehlt in der Bank (Platzhalter: Kettenname) – Nr. 1
%% TODO Anweisung: ein Satz über der ersten Teilaufgabe fehlt in der Bank – Nr. 1
\begin{aufgabe}{Quadrieren auch negativer Zahlen und Dezimalzahlen, Quadrat vor Punkt vor Strich: (−7)² = 49, nicht −49; 0,5² = 0,25}
\begin{teile}
\teil $6^2$ – wie viel? \leerfeld
\teil $1{,}2^2$ – wie viel? \leerfeld
\end{teile}
\end{aufgabe}

%% TODO Titel: Ich-kann-Satz fehlt in der Bank (Platzhalter: Kettenname) – Nr. 2
%% TODO Anweisung: ein Satz über der ersten Teilaufgabe fehlt in der Bank – Nr. 2
\begin{aufgabe}{Koordinaten lesen und eintragen, vier Quadranten, (x | y)-Reihenfolge, Kästchenraster mit 0,5-Einteilung}
\begin{teile}
\teil Punkt $A$ – Koordinaten?

\begin{ksys}[xmin=-4,xmax=4,ymin=-4,ymax=4,ablesen]
\punkt{3}{2}{A}
\end{ksys}
A( \leerfeld | \leerfeld )
\teil Punkt $C$ – Koordinaten?

\begin{ksys}[xmin=-5,xmax=4,ymin=-4,ymax=4,ablesen]
\punkt{-4}{-3}{C}
\end{ksys}
C( \leerfeld | \leerfeld )
\end{teile}
\end{aufgabe}

%% TODO Titel: Ich-kann-Satz fehlt in der Bank (Platzhalter: Kettenname) – Nr. 3
%% TODO Anweisung: ein Satz über der ersten Teilaufgabe fehlt in der Bank – Nr. 3
\begin{aufgabe}{Lineare Funktion f(x) = m·x + n: Gerade zeichnen, Funktionswert, Punktprobe, Nullstelle}
\begin{teile}
\teil $g(x) = 2x + 1$ – $g(4)$? \leerfeld
\teil $g(x) = -3x + 2$ – $g(-2)$? \leerfeld
\end{teile}
\end{aufgabe}

%% TODO Titel: Ich-kann-Satz fehlt in der Bank (Platzhalter: Kettenname) – Nr. 4
%% TODO Anweisung: ein Satz über der ersten Teilaufgabe fehlt in der Bank – Nr. 4
\begin{aufgabe}{Binomische Formel ausmultiplizieren und zusammenfassen, auch mit Minus in der Klammer}
\begin{teile}
\teil $(x + 4)^2$ – ausmultipliziert? \leerfeld
\teil $(x - 6)^2$ – ausmultipliziert? \leerfeld
\end{teile}
\end{aufgabe}

%% TODO Titel: Ich-kann-Satz fehlt in der Bank (Platzhalter: Kettenname) – Nr. 5
%% TODO Anweisung: ein Satz über der ersten Teilaufgabe fehlt in der Bank – Nr. 5
\begin{aufgabe}{Lineare Gleichung lösen, Terme mit x auf eine Seite bringen, Schreibform mit Strich}
\begin{gleichungsraster}[2]
\gl{\text{$3x + 2 = 14$ – $x$?}} &
\gl{\text{$2x + 9 = -3$ – $x$?}} \\
\end{gleichungsraster}
\end{aufgabe}

%% TODO Titel: Ich-kann-Satz fehlt in der Bank (Platzhalter: Kettenname) – Nr. 6
%% TODO Anweisung: ein Satz über der ersten Teilaufgabe fehlt in der Bank – Nr. 6
\begin{aufgabe}{Schnittpunkt zweier Geraden durch Gleichsetzen}
\begin{gleichungsraster}[2]
\gl{\text{$g(x) = x + 1$ und $h(x) = -x + 5$ – Schnittpunkt?}} &
\gl{\text{$g(x) = 3x + 4$ und $h(x) = x - 2$ – Schnittpunkt?}} \\
\end{gleichungsraster}
\end{aufgabe}

%% TODO Titel: Ich-kann-Satz fehlt in der Bank (Platzhalter: Kettenname) – Nr. 7
%% TODO Anweisung: ein Satz über der ersten Teilaufgabe fehlt in der Bank – Nr. 7
\begin{aufgabe}{Quadratwurzel ziehen, auch mit dem Taschenrechner, Näherungswert runden; aus einer negativen Zahl gibt es keine Wurzel}
\begin{teile}
\teil $\sqrt{64}$ – wie viel? \leerfeld
\teil $\sqrt{7}$ – auf zwei Stellen gerundet? \leerfeld
\end{teile}
\end{aufgabe}

%% TODO Titel: Ich-kann-Satz fehlt in der Bank (Platzhalter: Kettenname) – Nr. 8
%% TODO Anweisung: ein Satz über der ersten Teilaufgabe fehlt in der Bank – Nr. 8
\begin{aufgabe}{Quadratische Gleichung durch Wurzelziehen lösen und Normalform mit der p-q-Formel lösen, beide Lösungen angeben}
\begin{gleichungsraster}[2]
\gl{\text{$x^2 = 36$ – Lösungen?}} &
\gl{\text{$x^2 + 4x - 5 = 0$ – Lösungen?}} \\
\end{gleichungsraster}
\end{aufgabe}

%% TODO Titel: Ich-kann-Satz fehlt in der Bank (Platzhalter: Kettenname) – Nr. 9
%% TODO Anweisung: ein Satz über der ersten Teilaufgabe fehlt in der Bank – Nr. 9
\begin{aufgabe}{Binomische Formel ausmultiplizieren und zusammenfassen, auch mit Minus in der Klammer – Fehler finden}
\begin{teile}
\teil Tim rechnet so: \rechnung{(x - 5)^2 &= x^2 - 25} Finde den Fehler und rechne richtig. \leerfeld
\end{teile}
\end{aufgabe}

%% TODO Titel: Ich-kann-Satz fehlt in der Bank (Platzhalter: Kettenname) – Nr. 10
%% TODO Anweisung: ein Satz über der ersten Teilaufgabe fehlt in der Bank – Nr. 10
\begin{aufgabe}{Binomische Formel ausmultiplizieren und zusammenfassen, auch mit Minus in der Klammer – selbst rechnen}
\begin{teile}
\teil $(x + 7)^2$ – ausmultipliziert? \leerfeld
\end{teile}
\end{aufgabe}

\clearpage
% Einheit 1 von 4 – Normalparabel und Streckfaktor (bank/quadratische-funktionen/e1.jsonl, zusammenbau v0.1)
%% TODO Zweigzeile Teil 1: was hier gelernt wird (Satz in Schülersprache aus dem Einheitstitel) – Einheit 1
\einheitenkopf[e1]{Einheit 1 von 4 · Normalparabel und Streckfaktor}
\zweigzeile{ab Kl. 9, je nach Buch bis Kl. 10 · P10 oft}
\setcounter{aufgabe}{10}

%% TODO Titel: Ich-kann-Satz fehlt in der Bank (Platzhalter: Kettenname) – Nr. 11
%% TODO Anweisung: ein Satz über der ersten Teilaufgabe fehlt in der Bank – Nr. 11
\begin{aufgabe}{Normalparabel und Streckfaktor}
%% TODO \rechenplatz{n}: Schrittzahl des Grundfalls fehlt in der Bank (nur bei fester Schreibform, 2.3 b)
\begin{teile}
\teil $f(x) = 4x - 1$ – Gerade oder Parabel? \kreuz{Gerade} \kreuz{Parabel}
\teil Wertetabelle zu $f(x) = x^2$ – fülle aus.
\wertetabelle{x}{f(x)}{-2,-1,0,1,2}
\teil Trage die Punkte ein und zeichne die Parabel zu $f(x) = x^2$ als Bogen.
\wertetabelle[9,4,1,0,1,4,9]{x}{f(x)}{-3,-2,-1,0,1,2,3}

\begin{ksys}[xmin=-4,xmax=4,ymin=-1,ymax=10]
\end{ksys}
\teil Normalparabel $f(x) = x^2$ – Scheitel und Symmetrieachse?

\begin{ksys}[xmin=-4,xmax=4,ymin=-1,ymax=9,ablesen]
\parabel{1}{0}{0}{f}
\end{ksys}
S( \leerfeld | \leerfeld ), Symmetrieachse: \leerfeld
\teil $f(x) = x^2$ – $f(-6)$? f(-6) = \leerfeld
\teil Wertetabelle zu $f(x) = 4x^2$ – fülle aus.
\wertetabelle{x}{f(x)}{-2,-1,0,1,2}
\teil $f(x) = -4x^2$ – \kreuz{nach oben} \kreuz{nach unten} \\ \kreuz{schmaler} \kreuz{breiter} als die Normalparabel
\teil Welche Tabelle gehört zu $f(x) = 0{,}5x^2$? \\ A: \wertetabelle[2,0{,}5,0,0{,}5,2]{x}{y}{-2,-1,0,1,2} \\ B: \wertetabelle[1,0{,}25,0,0{,}25,1]{x}{y}{-2,-1,0,1,2} \\ C: \wertetabelle[4,1,0,1,4]{x}{y}{-2,-1,0,1,2} \\ \kreuz{A} \kreuz{B} \kreuz{C}
\teil $f(x) = 0{,}5x^2$ – zeichne die Parabel mit den halbierten Werten der Normalparabel.

\begin{ksys}[xmin=-5,xmax=5,ymin=-1,ymax=9]
\end{ksys}
\teil Welche Gleichung gehört zum Graphen? \\ \kreuz{$f(x) = x^2 + 3$} \\ \kreuz{$f(x) = -x^2 - 3$} \\ \kreuz{$f(x) = x^2 - 3$} \\ \kreuz{$f(x) = -x^2 + 3$}

\begin{ksys}[xmin=-4,xmax=4,ymin=-5,ymax=5,ablesen]
\funktion{-1*\x^2+(3)}{f}
\end{ksys}
\teil Welche Wertetabelle gehört zu $y = 4x^2$? \\ A: \wertetabelle[-8,-4,0,4,8]{x}{y}{-2,-1,0,1,2} \\ B: \wertetabelle[16,4,0,4,16]{x}{y}{-2,-1,0,1,2} \\ C: \wertetabelle[64,16,0,16,64]{x}{y}{-2,-1,0,1,2} \\ \kreuz{A} \kreuz{B} \kreuz{C} \hfill (P10 2026 FOR)
\teil Ein Stein fällt von einer $45\,$m hohen Brücke. Der Graph zeigt seine Höhe $h$ (in m) nach $t$ Sekunden. Welche Gleichung passt? Kreuze an und begründe. \\ \kreuz{$h(t) = 5t^2 + 45$} \\ \kreuz{$h(t) = 5t^2 - 45$} \\ \kreuz{$h(t) = -5t^2 + 45$} \\ \kreuz{$h(t) = -5t^2 - 45$} \hfill (P10 2015 OS)

\begin{ksys}[xmin=0,xmax=4,ymin=0,ymax=50,xstep=0.5,ystep=5,xlabel=$t$ in s,ylabel=$h$ in m,ablesen]
\funktionab{-5*\x^2+45}{h}{0}{3}
\end{ksys}
\end{teile}
\end{aufgabe}

%% TODO Titel: Ich-kann-Satz fehlt in der Bank (Platzhalter: Kettenname) – Nr. 12
%% TODO Anweisung: ein Satz über der ersten Teilaufgabe fehlt in der Bank – Nr. 12
\begin{aufgabe}{Normalparabel und Streckfaktor – Fehler finden, Begründen, Darstellungswechsel, Anwendung}
\begin{teile}
\teil Mia füllt die Tabelle zu $f(x) = x^2$ so aus: \\ \wertetabelle[-16,-4,0,4,16]{x}{f(x)}{-4,-2,0,2,4} \\ Finde den Fehler und korrigiere.
\teil $f(x) = x^2$ hat bei $x = 8$ und bei $x = -8$ denselben Wert. Warum?
\teil Die Tabelle gehört zu $f(x) = a \cdot x^2$. \\ \wertetabelle[7,28,63]{x}{f(x)}{1,2,3} \\ Wie heißt $a$? a = \leerfeld
\teil Der Bogen einer Brücke folgt $h(x) = -0{,}02x^2$ ($x$ und $h$ in m, der Ursprung liegt im höchsten Punkt des Bogens). Wie tief liegt der Bogen $20\,$m neben der Mitte unter dem höchsten Punkt?

\begin{ksys}[xmin=-30,xmax=30,ymin=-20,ymax=5,xstep=10,ystep=5,xlabel=$x$ in m,ylabel=$h$ in m]
\funktion{-0.02*\x^2}{h}
\end{ksys}
\leerfeld[m]
\end{teile}
\end{aufgabe}

\clearpage
% Einheit 2 von 4 – Scheitelpunktform (bank/quadratische-funktionen/e2.jsonl, zusammenbau v0.1)
%% TODO Zweigzeile Teil 1: was hier gelernt wird (Satz in Schülersprache aus dem Einheitstitel) – Einheit 2
\einheitenkopf[e2]{Einheit 2 von 4 · Scheitelpunktform}
\zweigzeile{ab Kl. 9, je nach Buch bis Kl. 10 · P10 oft}
\setcounter{aufgabe}{12}

%% TODO Titel: Ich-kann-Satz fehlt in der Bank (Platzhalter: Kettenname) – Nr. 13
%% TODO Anweisung: ein Satz über der ersten Teilaufgabe fehlt in der Bank – Nr. 13
\begin{aufgabe}{Scheitelpunktform}
%% TODO \rechenplatz{n}: Schrittzahl des Grundfalls fehlt in der Bank (nur bei fester Schreibform, 2.3 b)
\begin{teile}
\teil $f(x) = (x - 4)^2 + 2$ – kreise die Zahl in der Klammer ein. Bei welchem $x$ liegt der Scheitel? x = \leerfeld
\teil $f(x) = (x - 2)^2 + 5$ – Scheitel? S( \leerfeld | \leerfeld )
\teil $f(x) = (x + 4)^2 + 2$ – Scheitel? S( \leerfeld | \leerfeld )
\teil Lies den Scheitel der Parabel ab.

\begin{ksys}[xmin=-5,xmax=1,ymin=-1,ymax=8,ablesen]
\parabel{1}{-2}{3}{f}
\end{ksys}
S( \leerfeld | \leerfeld )
\teil $f(x) = (x - 2)^2 + 1$ – skizziere die Parabel.

\begin{ksys}[xmin=-2,xmax=6,ymin=-1,ymax=10]
\end{ksys}
\teil Nach oben geöffnete Normalparabel mit Scheitel $S(5|2)$ – Gleichung? f(x) = \leerfeld
\teil Gleichung der Parabel?

\begin{ksys}[xmin=-1,xmax=5,ymin=-4,ymax=2,ablesen]
\parabel{1}{2}{-3}{f}
\end{ksys}
f(x) = \leerfeld
\teil $f(x) = (x - 1)^2 + 2$ wird um $4$ nach unten verschoben – Gleichung der neuen Parabel $g$? g(x) = \leerfeld
\teil $f(x) = (x - 5)^2 + 1$ wird an der $x$-Achse gespiegelt – Gleichung des Spiegelbilds $g$? g(x) = \leerfeld
\teil $f(x) = (x - 4)^2 + 2$ wird an der $y$-Achse gespiegelt – Gleichung des Spiegelbilds $g$? g(x) = \leerfeld
\teil Gib eine nach oben geöffnete Normalparabel an, deren Scheitel auf der $x$-Achse rechts vom Ursprung liegt. f(x) = \leerfeld
\teil $f(x) = (x - 1)^2 + 2$ und $g(x) = (x - 1)^2 - 4$ – wie liegen die Parabeln zueinander? Haben sie gemeinsame Punkte? Begründe ohne Rechnung.
\teil $f(x) = 2(x - 4)^2 + 1$ – Scheitel, Öffnung, schmaler oder breiter als die Normalparabel? Punkt eine Einheit rechts vom Scheitel? S( \leerfeld | \leerfeld ), Punkt ( \leerfeld | \leerfeld )
\teil $f(x) = (x - 4)^2 - 5$ wird erst um $5$ nach oben verschoben und dann an der $x$-Achse gespiegelt. Gleichung der neuen Parabel $g$? \hfill (P10 2017 OS) \\ g(x) = \leerfeld
\teil Die Parabel $f(x) = (x + 4)^2 - 2$ wird an der $y$-Achse gespiegelt. Skizziere das Spiegelbild $g$ und gib seine Gleichung an. \hfill (P10 2020 OS) \\

\begin{ksys}[xmin=-7,xmax=7,ymin=-3,ymax=7]
\parabel{1}{-4}{-2}{f}
\end{ksys}
g(x) = \leerfeld
\teil Skizziere eine nach unten geöffnete Normalparabel $f$ mit Scheitel auf der $y$-Achse, die keine Nullstelle hat. Gib ihre Gleichung und die Gleichung ihres Spiegelbilds $g$ an der $x$-Achse an. \hfill (P10 2018 OS) \\

\begin{ksys}[xmin=-4,xmax=4,ymin=-10,ymax=10,ystep=2]
\end{ksys}
f(x) = \leerfeld , g(x) = \leerfeld
\end{teile}
\end{aufgabe}

%% TODO Titel: Ich-kann-Satz fehlt in der Bank (Platzhalter: Kettenname) – Nr. 14
%% TODO Anweisung: ein Satz über der ersten Teilaufgabe fehlt in der Bank – Nr. 14
\begin{aufgabe}{Scheitelpunktform – Fehler finden, Begründen, Darstellungswechsel, Anwendung}
\begin{teile}
\teil Ben liest bei $f(x) = (x + 6)^2 - 2$ den Scheitel $S(6|{-2})$ ab. Finde den Fehler und korrigiere.
\teil Warum hat $f(x) = (x - 4)^2 + 1$ bei $x = 4$ ihren kleinsten Wert?
\teil Wahr oder falsch? \\ Die Gleichung ist $f(x) = (x + 1)^2 - 4$. \kreuz{wahr} \kreuz{falsch} \\ Der Scheitel ist $S(1|{-4})$. \kreuz{wahr} \kreuz{falsch} \\ Die Parabel ist nach unten geöffnet. \kreuz{wahr} \kreuz{falsch}

\begin{ksys}[xmin=-2,xmax=4,ymin=-5,ymax=1,ablesen]
\parabel{1}{1}{-4}{f}
\end{ksys}
\teil Beim Kugelstoßen gilt für die Flugbahn $h(x) = -0{,}05(x - 6)^2 + 3{,}8$ ($x$: Entfernung in m, $h$: Höhe in m). Wie hoch fliegt die Kugel höchstens, und in welcher Entfernung?

\begin{ksys}[xmin=0,xmax=16,ymin=0,ymax=5,xstep=2,xlabel=$x$ in m,ylabel=$h$ in m]
\funktionab{-0.05*(\x-6)^2+3.8}{h}{0}{14.7}
\end{ksys}
\leerfeld[m] hoch, \leerfeld[m] entfernt
\end{teile}
\end{aufgabe}

\clearpage
% Einheit 3 von 4 – Normalform (bank/quadratische-funktionen/e3.jsonl, zusammenbau v0.1)
%% TODO Zweigzeile Teil 1: was hier gelernt wird (Satz in Schülersprache aus dem Einheitstitel) – Einheit 3
\einheitenkopf[e3]{Einheit 3 von 4 · Normalform}
\zweigzeile{ab Kl. 9, je nach Buch bis Kl. 10 · P10 oft}
\setcounter{aufgabe}{14}

%% TODO Titel: Ich-kann-Satz fehlt in der Bank (Platzhalter: Kettenname) – Nr. 15
%% TODO Anweisung: ein Satz über der ersten Teilaufgabe fehlt in der Bank – Nr. 15
\begin{aufgabe}{Normalform}
%% TODO \rechenplatz{n}: Schrittzahl des Grundfalls fehlt in der Bank (nur bei fester Schreibform, 2.3 b)
\begin{teile}
\teil $f(x) = x^2 + 4x - 1$ – welche Form? \kreuz{Scheitelpunktform} \kreuz{Normalform}
\teil $f(x) = x^2 + 3x + 5$ – Schnittpunkt mit der $y$-Achse? ( \leerfeld | \leerfeld )
\teil $f(x) = x^2 + 4x - 2$ – $f(-3)$? f(-3) = \leerfeld
\end{teile}
\begin{gleichungsraster}[2]
\gl{\text{$f(x) = (x - 2)^2 + 5$ – Normalform?}} &
\gl{\text{$f(x) = (x + 4)^2 + 1$ – Normalform?}} \\
\gl{\text{Stimmt $(x + 5)^2 - 8 = x^2 + 10x + 17$? Weise nach oder widerlege.}} &
\end{gleichungsraster}
\begin{teile}
\teil $f(x) = x^2 - 4x + 1$ – lies den Scheitel am Graphen ab.

\begin{ksys}[xmin=-1,xmax=5,ymin=-4,ymax=2,ablesen]
\parabel{1}{2}{-3}{f}
\end{ksys}
S( \leerfeld | \leerfeld )
\teil $f(x) = x^2 - 6x + 5$ – lies den Scheitel ab, stelle die Scheitelpunktform auf und prüfe mit $x = 0$.

\begin{ksys}[xmin=-1,xmax=6,ymin=-5,ymax=1,ablesen]
\parabel{1}{3}{-4}{f}
\end{ksys}
f(x) = \leerfeld
\teil $f(x) = x^2 + 2x - 3$ – wahr oder falsch? \\ $(2|5)$ liegt auf $f$. \kreuz{wahr} \kreuz{falsch} \\ $f$ ist die um $1$ nach links und um $4$ nach unten verschobene Normalparabel. \kreuz{wahr} \kreuz{falsch} \\ Der Schnittpunkt mit der $y$-Achse ist $(-3|0)$. \kreuz{wahr} \kreuz{falsch}

\begin{ksys}[xmin=-4,xmax=3,ymin=-5,ymax=6,ablesen]
\parabel{1}{-1}{-4}{f}
\end{ksys}
\teil Die Abbildung zeigt die Parabel $f(x) = x^2 - 8x + 13$. Gib den Scheitel an und notiere die Scheitelpunktform. \hfill (P10 2025 OS) \\

\begin{ksys}[xmin=-1,xmax=7,ymin=-4,ymax=2,ablesen]
\parabel{1}{4}{-3}{f}
\end{ksys}
S( \leerfeld | \leerfeld ), f(x) = \leerfeld
\end{teile}
\begin{gleichungsraster}[2]
\gl{\text{Weise nach: $(x + 4)^2 - 3 = x^2 + 8x + 13$. (P10 2017 OS)}} & \\
\end{gleichungsraster}
\end{aufgabe}

%% TODO Titel: Ich-kann-Satz fehlt in der Bank (Platzhalter: Kettenname) – Nr. 16
%% TODO Anweisung: ein Satz über der ersten Teilaufgabe fehlt in der Bank – Nr. 16
\begin{aufgabe}{Normalform – Fehler finden, Begründen}
\begin{gleichungsraster}[2]
\gl{\text{Nina rechnet: \rechnung{(x - 4)^2 + 1 &= x^2 - 16 + 1 \\ &= x^2 - 15} Finde den Fehler und korrigiere.}} &
\end{gleichungsraster}
\begin{teile}
\teil Der Graph von $f(x) = x^2 - 7x + 4$ schneidet die $y$-Achse im Punkt $(0|4)$. Warum?
\end{teile}
\end{aufgabe}

\clearpage
% Einheit 4 von 4 – Nullstellen und Schnittpunkte berechnen (bank/quadratische-funktionen/e4.jsonl, zusammenbau v0.1)
%% TODO Zweigzeile Teil 1: was hier gelernt wird (Satz in Schülersprache aus dem Einheitstitel) – Einheit 4
\einheitenkopf[e4]{Einheit 4 von 4 · Nullstellen und Schnittpunkte berechnen}
\zweigzeile{ab Kl. 10 (am Gymnasium ab Kl. 9) · P10 oft}
\setcounter{aufgabe}{16}

%% TODO Titel: Ich-kann-Satz fehlt in der Bank (Platzhalter: Kettenname) – Nr. 17
%% TODO Anweisung: ein Satz über der ersten Teilaufgabe fehlt in der Bank – Nr. 17
\begin{aufgabe}{Nullstellen und Schnittpunkte}
%% TODO \rechenplatz{n}: Schrittzahl des Grundfalls fehlt in der Bank (nur bei fester Schreibform, 2.3 b)
\begin{teile}
\teil Parabel $f(x) = x^2 - 2x$ und Gerade $g(x) = x + 4$ – welche Gleichung setzt du an? Nur hinschreiben. \leerfeld = \leerfeld
\end{teile}
\begin{gleichungsraster}[2]
\gl{\text{$f(x) = (x - 1)^2 - 4$ – Nullstellen?}} &
\end{gleichungsraster}
\begin{teile}
\teil $f(x) = (x - 2)^2 + 3$ – wie viele Nullstellen? Begründe ohne Rechnung.
\end{teile}
\begin{gleichungsraster}[2]
\gl{\text{$f(x) = x^2 - 2x - 8$ – Nullstellen?}} &
\gl{\text{$f(x) = x^2 - 4x + 1$ – Nullstellen? Runde auf zwei Stellen.}} \\
\gl{\text{$f(x) = 3x^2 - 3x - 18$ – Nullstellen?}} &
\gl{\text{$f(x) = x^2 + x + 1$ – für welche $x$ ist $f(x) = 7$?}} \\
\gl{\text{$f(x) = x^2 + x - 4$ und $g(x) = 2x + 2$ – an welchen Stellen $x$ schneiden sich Parabel und Gerade?}} &
\gl{\text{$f(x) = (x - 2)^2 - 1$ und $g(x) = x - 3$ – an welchen Stellen $x$ schneiden sich Parabel und Gerade?}} \\
\gl{\text{$f(x) = x^2 - 4x + 1$ und $g(x) = 2x - 4$ – Schnittpunkte?}} &
\end{gleichungsraster}
\begin{teile}
\teil $f(x) = x^2 - 5$, $g(x) = 2x + 3$ – ist $P(4|11)$ ein Schnittpunkt? \janein
\teil $f(x) = (x - 2)^2 + 1$ – gib eine Gerade $g$ an, die mit der Parabel keinen Punkt gemeinsam hat. g(x) = \leerfeld
\end{teile}
\begin{gleichungsraster}[2]
\gl{\text{$f(x) = x^2 + 3x - 2$ und $h(x) = x^2 - x + 6$ – Schnittpunkt?}} &
\gl{\text{Die Parabel $f(x) = (x - 4)^2 - 6$ und die Gerade $g(x) = -2x + 5$ schneiden sich in zwei Punkten. Berechne ihre Koordinaten. (P10 2022 OS)}} \\
\gl{\text{Berechne die Schnittpunkte der Parabel $f(x) = x^2 - 10x + 22$ mit der $x$-Achse. Runde auf zwei Stellen. (P10 2025 OS)}} &
\gl{\text{Die Parabel $f(x) = -x^2 + 4x + 2$ ist nach unten geöffnet. Für welche $x$ ist $f(x) = -10$? (P10 2023 OS)}} \\
\end{gleichungsraster}
\end{aufgabe}

%% TODO Titel: Ich-kann-Satz fehlt in der Bank (Platzhalter: Kettenname) – Nr. 18
%% TODO Anweisung: ein Satz über der ersten Teilaufgabe fehlt in der Bank – Nr. 18
\begin{aufgabe}{Nullstellen und Schnittpunkte – Fehler finden, Begründen, Darstellungswechsel, Anwendung}
\begin{gleichungsraster}[2]
\gl{\text{Lina berechnet die Nullstellen von $f(x) = (x + 1)^2 - 25$: \rechnung{0 &= (x + 1)^2 - 25 \\ 25 &= (x + 1)^2 \\ x + 1 &= 5 \\ x &= 4} Finde den Fehler und korrigiere.}} &
\end{gleichungsraster}
\begin{teile}
\teil Warum liegt jede Nullstelle einer Parabel auf der $x$-Achse?
\teil $f(x) = (x - 2)^2 - 3$ und $g(x) = x - 3$ – lies die Schnittpunkte am Graphen ab und bestätige einen durch Einsetzen.

\begin{ksys}[xmin=-1,xmax=5,ymin=-4,ymax=2,ablesen]
\parabel{1}{2}{-3}{f}
\gerade{1}{-3}{g}
\end{ksys}
$S_1($ \leerfeld | \leerfeld ), $S_2($ \leerfeld | \leerfeld )
\end{teile}
\begin{gleichungsraster}[2]
\gl{\text{Ein Ball fliegt auf der Bahn $h(x) = -0{,}1(x - 6)^2 + 4{,}9$ ($x$: Entfernung in m, $h$: Höhe in m). Wie weit fliegt der Ball, bis er auf dem Boden aufkommt?}} & \\
\end{gleichungsraster}
\end{aufgabe}

% Abhakseite „Das kann ich“ – Zone nur im Gesamt (\mitzone)
%% TODO Ich-kann-Sätze: Platzhalter wie in den Aufgabendateien
\begin{abhakseite}
\ifdefined\mitzone
\abhakgruppe{Kennst du schon}
\abhak{1}{Quadrieren auch negativer Zahlen und Dezimalzahlen, Quadrat vor Punkt vor Strich: (−7)² = 49, nicht −49; 0,5² = 0,25}
\abhak{2}{Koordinaten lesen und eintragen, vier Quadranten, (x | y)-Reihenfolge, Kästchenraster mit 0,5-Einteilung}
\abhak{3}{Lineare Funktion f(x) = m·x + n: Gerade zeichnen, Funktionswert, Punktprobe, Nullstelle}
\abhak{4}{Binomische Formel ausmultiplizieren und zusammenfassen, auch mit Minus in der Klammer}
\abhak{5}{Lineare Gleichung lösen, Terme mit x auf eine Seite bringen, Schreibform mit Strich}
\abhak{6}{Schnittpunkt zweier Geraden durch Gleichsetzen}
\abhak{7}{Quadratwurzel ziehen, auch mit dem Taschenrechner, Näherungswert runden; aus einer negativen Zahl gibt es keine Wurzel}
\abhak{8}{Quadratische Gleichung durch Wurzelziehen lösen und Normalform mit der p-q-Formel lösen, beide Lösungen angeben}
\abhak{9}{Binomische Formel ausmultiplizieren und zusammenfassen, auch mit Minus in der Klammer – Fehler finden}
\abhak{10}{Binomische Formel ausmultiplizieren und zusammenfassen, auch mit Minus in der Klammer – selbst rechnen}
\fi
\abhakgruppe{Einheit 1 · Normalparabel und Streckfaktor}
\abhak{11}{Normalparabel und Streckfaktor}
\abhak{12}{Normalparabel und Streckfaktor – Fehler finden, Begründen, Darstellungswechsel, Anwendung}
\abhakgruppe{Einheit 2 · Scheitelpunktform}
\abhak{13}{Scheitelpunktform}
\abhak{14}{Scheitelpunktform – Fehler finden, Begründen, Darstellungswechsel, Anwendung}
\abhakgruppe{Einheit 3 · Normalform}
\abhak{15}{Normalform}
\abhak{16}{Normalform – Fehler finden, Begründen}
\abhakgruppe{Einheit 4 · Nullstellen und Schnittpunkte berechnen}
\abhak{17}{Nullstellen und Schnittpunkte}
\abhak{18}{Nullstellen und Schnittpunkte – Fehler finden, Begründen, Darstellungswechsel, Anwendung}
\end{abhakseite}

\end{document}
